% Generated draft: requires mathematical and source review.
\documentclass[11pt]{article}
\usepackage[margin=25mm]{geometry}
\usepackage{fontspec}
\setmainfont{Latin Modern Roman}
\usepackage{amsmath,amssymb,mathtools,unicode-math}
\setmathfont{Latin Modern Math}
\usepackage{enumitem,needspace,xurl,hyperref}
\setlength{\parindent}{0pt}
\setlength{\parskip}{5pt}
\setlength{\emergencystretch}{3em}
\newcommand{\N}{\mathbb{N}}
\newcommand{\Z}{\mathbb{Z}}
\newcommand{\Q}{\mathbb{Q}}
\newcommand{\R}{\mathbb{R}}
\newcommand{\C}{\mathbb{C}}
\newcommand{\F}{\mathbb{F}}
\newcommand{\A}{\mathbb{A}}
\newcommand{\PP}{\mathbb{P}}
\newcommand{\E}{\mathbb{E}}
\newcommand{\eps}{\varepsilon}
\newcommand{\dd}{\,\mathrm{d}}
\begin{document}
\section*{UnsolvedMath ID: 30003459}
\textbf{Source problem code:} catalogue entry
\subsection{Definitions and mathematical statement}
For a representation \(V\), an honest subrepresentation is an invariant subspace of \(V\) with its induced action; equivalently, a representation \(W\) occurs as a subrepresentation if there is an injective equivariant map \(W\hookrightarrow V\). A subquotient is a representation isomorphic to \(U/U'\), where \(U'\subseteq U\subseteq V\) are invariant subspaces. In a semisimple direct-sum decomposition indexed by irreducible isomorphism classes, the multiplicity of a class is the number of summands belonging to that class.

The supplied record names a cuspidal automorphic representation space \(A_{G,\mathrm{cusp},\chi}\) and an indexing collection \(\Pi_{G,\mathrm{cusp},\chi}\), and displays
\[
A_{G,\mathrm{cusp},\chi}
=\bigoplus_{\pi\in\Pi_{G,\mathrm{cusp},\chi}}\pi.
\]
Its requested assertion is that every \(\pi\in\Pi_{G,\mathrm{cusp},\chi}\) admits an equivariant embedding
\[
\iota_\pi:\pi\hookrightarrow A_{G,\mathrm{cusp},\chi},
\]
and that each irreducible isomorphism class represented in this collection occurs with multiplicity exactly one in the displayed decomposition. Multiplicity one concerns the number of copies of an isomorphism class, not uniqueness of the embedding map.

\textbf{Unresolved scope.} The record does not specify the underlying global field, the group \(G\), the meaning or admissibility conditions of \(\chi\), the acting group, the coefficient field, the representation category, or the precise definitions of \(A_{G,\mathrm{cusp},\chi}\) and \(\Pi_{G,\mathrm{cusp},\chi}\). It also does not specify whether the displayed sum is algebraic or a completed Hilbert direct sum, or whether its indices are distinct isomorphism classes. The assertion above therefore records the intended content schematically; it is not a universally quantified assertion for arbitrary groups or automorphic settings.

\textbf{Supplied status qualification.} The registry labels the item open, but its supplied statement-verification note says that the report states and proves the multiplicity-one assertion in a surrounding theorem, and describes the extraction as an answered expository question. Its statement status is \texttt{unrecoverable}. No complete unresolved theorem is identified by the supplied record.
\subsection{Short English statement}
The extracted question asks whether the indicated cuspidal automorphic representations occur as actual invariant subspaces, each with multiplicity one. Essential hypotheses and the meaning of the decomposition are missing. Although the registry labels the item open, its supplied verification note describes it as an answered expository question; this is a statement-only draft requiring human review.
\subsection{Sources}
\begin{itemize}
\item Local UnsolvedMath registry, \texttt{data/unsolvedmath/problems.json}, permanent ID \(30003459\), source problem code \texttt{catalogue entry}, title \emph{Multiplicity One in Cuspidal Automorphic Decompositions}. No individual proposer is supplied; the proposed year is recorded as 2017.
\item Supplied report citation: \emph{Arbeitsgemeinschaft: Higher Gross Zagier Formulas (2018)}, Oberwolfach Reports, pp.~1067--1134, DOI \texttt{10.4171/owr/2017/18}. Supplied source and literature-source link, labeled ``Original Oberwolfach report'': \url{https://doi.org/10.4171/owr/2017/18}. The record attributes the extraction to the report's workshop discussions; no exact theorem or passage locator is supplied.
\item Literature assessment checked at \texttt{2026-08-21}: supplied dated curation, not independently verified. It classifies the item as open because the extraction asks for the content of a displayed definition-and-corollary passage and its intended unresolved target is unclear. The background requests independent verification of time-sensitive status and cited claims.
\item The separately supplied statement-verification note says that the official report states and proves the multiplicity-one assertion in the surrounding theorem, identifies the extraction as an answered expository question, and says no separate clean statement is promoted. These are supplied source-verification claims, not independent verification in this draft.
\end{itemize}
\end{document}
